\begin{exsol}{20}{MWG Exercise 8.B.4}{262}{Advanced}{B3}
\asked Show that the order of deletion does not matter for the set of strategies that survive
iterated deletion of strictly dominated strategies.

\method Proof for finite games with dominance by pure strategies (MWG \S8.B, pp.~239--240;
Lecture~1, slide~29). Two facts do the work: strict dominance is inherited by smaller opponent
sets, and in a finite game a strictly dominated strategy has an undominated dominator.

\solution
\textbf{Setup.} $G = (I, \{S_i\}, \{u_i\})$ with every $S_i$ finite. For a product set
$T = \prod_i T_i \subseteq S$, say $s_i \in T_i$ is \emph{dominated in $T$} if some $t_i \in T_i$
satisfies $u_i(t_i, \sm) > u_i(s_i, \sm)$ for all $\sm \in T_{-i}$. A \emph{deletion sequence} is
$S = T^0 \supsetneq T^1 \supsetneq \dots \supsetneq T^K$, where each step removes a nonempty set
of strategies that are dominated in the current set (any players, any number, in any order), and
where nothing is dominated in $T^K$. The lecture's procedure (remove everything dominated each
round) is one such sequence. \emph{Claim:} every deletion sequence ends at the same $T^K$.

\textbf{Two facts.}
\begin{enumerate}[label=(\roman*)]
  \item \emph{Inheritance.} If $t_i$ beats $s_i$ against every $\sm \in A_{-i}$ and
  $B_{-i} \subseteq A_{-i}$, it beats $s_i$ against every $\sm \in B_{-i}$.
  \item \emph{Transitivity.} If $t_i$ beats $s_i$ and $r_i$ beats $t_i$ on the same set, then
  $r_i$ beats $s_i$ there.
\end{enumerate}

\textbf{Proof.} Let $(A^k)_{k=0}^{K}$ and $(B^m)_{m=0}^{M}$ be two deletion sequences with end
sets $A = A^K$ and $B = B^M$. We show $B \subseteq A^k$ for all $k$, by induction. It holds for
$k=0$. Suppose $B \subseteq A^k$ and that some $s_i \in B_i$ is deleted at step $k+1$ of the
$A$-sequence. Then some $t_i \in A_i^k$ beats $s_i$ on $A_{-i}^k \supseteq B_{-i}$, so by (i)
the set
\[
  D = \{\, r_i \in S_i : u_i(r_i, \sm) > u_i(s_i, \sm) \ \forall \sm \in B_{-i} \,\}
\]
is nonempty. We show that $D \cap B_i \neq \emptyset$. Then $s_i$ is dominated in $B$, which
contradicts the fact that the $B$-sequence stopped.

Suppose instead that every $r_i \in D$ is deleted at some step of the $B$-sequence. Let $m$ be the
last step at which an element of $D$ is deleted ($D$ is finite), and let $D_m$ be the elements of
$D$ deleted at step $m$. ``Beats on $B_{-i}^{m-1}$'' is a strict partial order on $D_m$
(irreflexive and transitive), so $D_m$ has an element $r^*$ that no other element of $D_m$
beats on $B_{-i}^{m-1}$. Since $r^*$ is deleted at step $m$, some $q \in B_i^{m-1}$ beats $r^*$
on $B_{-i}^{m-1} \supseteq B_{-i}$. By (i) and (ii), $q$ beats $s_i$ on $B_{-i}$, so $q \in D$.
Now $q \in B_i^{m-1}$ and $q \in D$, so either $q$ is deleted at step $m$, and then $q \in D_m$
beats $r^*$, contradicting the choice of $r^*$; or $q$ is deleted after $m$, contradicting the
choice of $m$; or $q$ is never deleted, so $q \in D \cap B_i$, contradicting our supposition.
Hence $D \cap B_i \ne \emptyset$.

So no strategy of $B$ is ever deleted by the $A$-sequence: $B \subseteq A$. Exchanging the roles
of the two sequences gives $A \subseteq B$, so $A = B$. \hfill$\square$

\textbf{Remarks.}
\begin{itemize}
  \item The selection of $r^*$ handles a subtle point: the strategy that dominates $s_i$ may itself
  be deleted, possibly in the same step. Finiteness guarantees that the chain of dominators
  ends at a strategy that survives.
  \item Finiteness matters. In infinite games a strategy can be dominated only by strategies that
  are themselves dominated, and different orders can then end at different sets. Dufwenberg and
  Stegeman (2002, \emph{Econometrica}) give such examples. In the continuous games of items~15
  and~21 there is no such problem, because each round's survivors are exactly the possible best
  replies.
  \item The argument breaks for \emph{weak} dominance at step (i): weak dominance on $A_{-i}$ need not
  survive on a subset $B_{-i}$, because the strict inequality may sit on a deleted $\sm$
  (item~19).
\end{itemize}

\interp The set $S^\infty$ is a property of the game, not of the analyst's bookkeeping. One may
delete in whatever order is convenient (one player at a time, one strategy at a time) and the
prediction is the same. This is part of what makes IESDS a well-defined solution concept.

\verify \texttt{check\_C.py} draws 3\,000 random finite games (between $2\times2$ and $5\times5$,
payoffs $0$--$6$). For each it runs the lecture's simultaneous procedure and five random
sequences that each delete a random nonempty subset of one random player's dominated
strategies. All end at the same set. \emph{Comparison with the MWG solutions manual}
(pp.~8-2--8-3): it shows that a strategy strictly dominated at some round will be deleted at the
next round, assuming that its dominator $s_i^*$ is ``not a strictly dominated strategy''. It does
not prove that such an undominated dominator exists (the finiteness step above), and it does not
complete the comparison of two whole sequences. The conclusion is the same; the proof above
fills those two gaps. The manual's note on weak dominance agrees with item~19.

\answer{In a finite game every deletion sequence ends at the same set. A strategy in one end set
$B$ can never be deleted by another sequence. If it were, a dominator chain, cut off at a maximal
element of the last-deleted batch, would produce a strategy in $B$ that dominates it on
$B_{-i}$, contradicting that $B$ is final. Strict dominance is inherited by smaller opponent
sets; weak dominance is not, which is why weak deletion is order-dependent.}
\end{exsol}
